\begin{table}[h!]
\caption{Minimum reporting checklist for evaluating link-prediction models on
biomedical interaction graphs. Each requirement targets one of the two failure modes
quantified in this paper; satisfying only one still permits a substantially inflated
headline number, since the two failure modes are independent (Results, Discussion).}
\label{tab:correct_protocol}
\small
\begin{tabular}{p{3.1cm}p{4.7cm}p{5.4cm}}
\hline
Requirement & Failure mode it prevents & How this study implements it \\
\hline
Leak-free edge split & Target-link inclusion: the encoder can recover a test edge's
label from its own neighborhood because that edge was never removed from message
passing (Figure~\ref{fig:edge_split}). & Held-out edges and their automatically-added
reverse relation are stripped from the message-passing graph before training; an
automated multi-check test verifies no held-out edge is reachable before any downstream
number is trusted. \\
Matched train/eval negative distribution & The ``mismatch trap'': evaluating a model
against a different negative distribution than it was trained on produces an artifact
of the mismatch, not a difficulty measurement (Figure~\ref{fig:neg_sampling_mechanism}).
& The same negative sampler (uniform-random or degree-matched) is enforced at both
training and evaluation time for every reported number. \\
Model-free baseline & Without a no-learning point of comparison, a model that has
learned a trivial shortcut (e.g.\ gene popularity) is indistinguishable from one that
has learned the task (Table~\ref{tab:model_free_baselines}). & At least one topological
heuristic (for link prediction) or simple non-graph classifier (for node
classification) is reported alongside every trained model, under the same protocol. \\
\hline
\end{tabular}
\end{table}
